\documentclass[UTF8,a4paper,11pt]{ctexart}
\usepackage{amsmath,amssymb}
\usepackage[margin=2.5cm]{geometry}
\usepackage[hidelinks]{hyperref}
\title{2023年全国乙卷文科第20题：切线与参数单调性}
\date{2026年10月24日}
\begin{document}
\maketitle
\noindent\textbf{固定来源：}\url{https://github.com/DxAThing/Gaokao-Math-Problems-Compilation/blob/f1ae169b58a75b3db6e35c670b3f0d5c78718188/content/2023/national_paper_b_liberal.tex#L246}

设 $x>0$，
\[
f_a(x)=\left(\frac1x+a\right)\ln(1+x).
\]

\paragraph{一、切线}
\[
f_a'(x)=-\frac{\ln(1+x)}{x^2}
          +\frac{1/x+a}{1+x}.
\]
当 $a=-1$ 时，$f_{-1}(1)=0$，$f_{-1}'(1)=-\ln2$，
故曲线在 $x=1$ 处的切线为
\[
\boxed{y=-(\ln2)(x-1)}.
\]

\paragraph{二、单调递增的必要充分条件}
令
\[
G(x)=\frac{(1+x)\ln(1+x)-x}{x^2},\qquad x>0.
\]
则 $f_a'(x)=[a-G(x)]/(1+x)$。再令
\[
H(x)=\frac{x^2}{2}+x-(1+x)\ln(1+x).
\]
由 $H(0)=0$ 与 $H'(x)=x-\ln(1+x)>0$，得
$H(x)>0$，即
\[
G(x)<\frac12\quad(x>0).
\]
另一方面，用两次洛必达法则得到
\[
\lim_{x\to0^+}G(x)
 =\lim_{x\to0^+}\frac{\ln(1+x)}{2x}
 =\frac12.
\]
故 $\sup_{x>0}G(x)=1/2$，且上确界不取到。
当 $a\ge1/2$ 时，$f_a'(x)>0$ 恒成立；
当 $a<1/2$ 时，取充分小的 $x>0$ 有
$G(x)>a$，于是导数为负，不符合题意。
所以
\[
\boxed{a\in\left[\frac12,+\infty\right)}.
\]

\bigskip
\noindent\textbf{审校状态：}仅独立数学核验；新源码尚未实际
XeLaTeX 编译。原题及第三方图示版权归上游相应权利主体。
\end{document}
